\SourceLeaf{247}%
magnitudinem elementi linearis in superficie curva exprimit, atque ante omnia significationem geometricam coëfficientium \(E,F,G\) examinemus. Iam in art. 5 monuimus, in superficie curva concipi posse duo systemata linearum, alterum, in quibus singulis sola \(p\) sit variabilis, \(q\) constans; alterum, in quibus sola \(q\) variabilis, \(p\) constans. Quodlibet punctum superficiei considerari potest tamquam intersectio lineae primi systematis cum linea secundi: tumque elementum lineae primae huic puncto adiacens et variationi \(\mathrm{d}p\) respondens erit \(=\BareRoot{E}\mathord{.}\mathrm{d}p\), nec non elementum lineae secundae respondens variationi \(\mathrm{d}q\) erit \(=\BareRoot{G}\mathord{.}\mathrm{d}q\); denique denotando per \(\omega\) angulum inter haec elementa, facile perspicitur, fieri \(\cos\omega=\frac{F}{\BareRoot{EG}}\). Area autem elementi parallelogrammatici in superficie curva inter duas lineas primi systematis, quibus respondent \(q,q+\mathrm{d}q\), atque duas lineas systematis secundi, quibus respondent \(p,p+\mathrm{d}p\), erit \(\BareRoot{(EG-FF)}\mathrm{d}p\mathord{.}\mathrm{d}q\).

Linea quaecunque in superficie curva ad neutrum illorum systematum pertinens, oritur, dum \(p\) et \(q\) concipiuntur esse functiones unius variabilis novae, vel altera illarum functio alterius. Sit \(s\) longitudo talis curvae ab initio arbitrario numerata et versus directionem utramvis pro positiva habita. Denotemus per \(\theta\) angulum, quem efficit elementum \(\mathrm{d}s=\BareRoot{(E\mathrm{d}p^2+2F\mathrm{d}p\mathord{.}\mathrm{d}q+G\mathrm{d}q^2)}\) cum linea primi systematis per initium elementi ducta, et quidem ne ulla ambiguitas remaneat, hunc angulum semper ab eo ramo illius lineae, in quo valores ipsius \(p\) crescunt, inchoari, et versus eam plagam positive accipi supponemus, versus quam valores ipsius \(q\) crescunt. His ita intellectis facile perspicitur haberi
\[
\begin{aligned}
\cos\theta\mathord{.}\mathrm{d}s&=\BareRoot{E}\mathord{.}\mathrm{d}p+\BareRoot{G}\mathord{.}\cos\omega\mathord{.}\mathrm{d}q=\frac{E\mathrm{d}p+F\mathrm{d}q}{\BareRoot{E}}\\
\sin\theta\mathord{.}\mathrm{d}s&=\BareRoot{G}\mathord{.}\sin\omega\mathord{.}\mathrm{d}q=\frac{\BareRoot{(EG-FF)}\mathord{.}\mathrm{d}q}{\BareRoot{E}}
\end{aligned}
\]
\Article{18}
Investigabimus nunc, quaenam sit conditio, ut haec linea sit brevissima. Quum ipsius longitudo \(s\) expressa sit per integrale
\[
s=\int\BareRoot{(E\mathrm{d}p^2+2F\mathrm{d}p\mathord{.}\mathrm{d}q+G\mathrm{d}q^2)}
\]
conditio minimi requirit, ut variatio huius integralis a mutatione infinite parva tractus lineae oriunda fiat \(=0\). Calculus ad propositum nostrum in hoc casu commodius absolvitur, si \(p\) tamquam functionem ipsius \(q\) consideramus. Quo
